\documentclass[12pt]{article}
\usepackage[utf8]{inputenc}
\usepackage[T1]{fontenc}
\usepackage{hyperref}
\title{Samosa Language Resources for K-12 Curriculum Development}
\author{Compiled by Claude Sonnet 5.5 (claude-sonnet-5-5) for Andr\'as Kornai}
\date{}
\begin{document}
\maketitle

\section{Introduction}
Samosa (ISO 639-3 swm, Glottocode samo1306) is a severely endangered Trans-New Guinea (Madang) language of Madang Province, Papua New Guinea. Wikipedia gives about 90 speakers (a 2000 figure), while Joshua Project lists about 200 people. Endangered Languages places it about 30 km inland from Alexishafen, near the headwaters of a tributary of the Gogol River. The searches found no school materials, grammar or dictionary. The name is easily confused with the Samo language of Burkina Faso, with Sam, and with Pamosu, which are different languages.

\section{Primary Educational Resources}
None found. The language is listed as having no written tradition.

\section{Linguistic and Reference Materials}
\begin{itemize}
\item Wikipedia article (a stub) with classification and speaker figures: \url{https://en.wikipedia.org/wiki/Samosa_language}
\item Endangered Languages Project entry, with location and pressure from larger languages such as Tok Pisin: \url{https://www.endangeredlanguages.com/lang/2648}
\end{itemize}

\section{Religious and Community Resources}
\begin{itemize}
\item Joshua Project profile of the Samosa people: \url{https://joshuaproject.net/people_groups/14678}
\end{itemize}

\section{Digital and Audio Resources}
No audio, video or corpus resources were found in the search results.

\section{Curriculum Development Context}
With only a few dozen to a couple of hundred speakers, mostly adults, and no orthography or published materials, any curriculum would begin with documentation: recording elders, agreeing on an orthography with the community, and producing simple bilingual materials with Tok Pisin and English. Nothing found suggests current school use.

\end{document}